\chapterspace{-2}
\chapter{Descenso fielmente plano}
\label{VIII}
\marginpar{195}

\section{Descenso de módulos cuasicoherentes}
\label{VIII.1}

Sea $\Sch$
\label{indnot:hb}\oldindexnot{$\Sch$|hyperpage}%
la categoría de los preesquemas. Procediendo como en~\Ref{VI.11}.b,
se obtiene que la categoría de los pares~$(X,F)$ formados por un
preesquema~$X$ y un módulo~$F$ sobre~$X$ (donde los morfismos se definen,
como en \loccit, mediante la noción de imagen directa de un módulo por un
morfismo de espacios anillados) puede considerarse como una categoría
fibrada sobre~$\Sch$; el funtor de cambio de base relativo a un morfismo
$f\colon X\to Y$ de~$\Sch$ es el funtor de imagen inversa de módulos
por~$f$. (Obsérvese que la categoría fibra en~$X\in\Ob(\Sch)$ de la
categoría fibrada anterior es la categoría \emph{opuesta} a la categoría
de los módulos sobre~$X$). Como la imagen inversa de un módulo
cuasicoherente es cuasicoherente, se ve que la subcategoría plena de la
categoría de los pares~$(X,F)$ formada por aquellos para los que $F$ es
cuasicoherente es una subcategoría fibrada de la categoría fibrada
anterior. (En cambio, si no se imponen hipótesis sobre~$f$, la imagen
directa de un módulo cuasicoherente no es, en general, un módulo
cuasicoherente). Llamaremos simplemente a esta categoría fibrada la
\emph{categoría fibrada de los módulos cuasicoherentes sobre los
preesquemas}.

Recordemos, por otra parte, que un morfismo~$f\colon X\to Y$ de espacios
anillados se llama \emph{fielmente plano}
\index{fielmente plano (morfismo)|hyperpage}%
\index{plano (morfismo fielmente)|hyperpage}%
si es \emph{plano} (\ie si, para todo~$x\in X$,
$\mathcal{O}_{X,x}$ es un módulo plano sobre~$\mathcal{O}_{Y,f(x)}$,
\eqref{IV}) y \emph{sobreyectivo}. Se dice que $f$ es un morfismo
\emph{cuasicompacto} si la imagen inversa por~$f$ de toda parte
cuasicompacta es cuasicompacta; cuando $f$ es un morfismo de
preesquemas, esto significa también que la imagen inversa por~$f$ de un
abierto afín de~$Y$ es reunión \emph{finita} de abiertos afines de~$X$.

\begin{theoreme}
\label{VIII.1.1}
Sea
\marginpar{196}
$\cal{F}$ la categoría fibrada de los módulos cuasicoherentes sobre los
preesquemas. Sea
\ifthenelse{\boolean{orig}}
{$g\colon S''\to S$}
{$g\colon S'\to S$}
un morfismo de preesquemas fielmente plano y cuasicompacto. Entonces $g$
es un morfismo de descenso efectivo para~$\cal{F}$.
\index{descenso efectivo (morfismo de)|hyperpage}%
\end{theoreme}

Recordemos\footnote{Admitiremos aquí la teoría general del descenso
\ifthenelse{\boolean{orig}}
{expuesta expuesta con detalle}
{expuesta con detalle}
en el artículo de J\ptbl\ \textsc{Giraud} citado en la nota
\eqref{footnotegiraud} al pie de la página del Aviso, trabajo que en lo
sucesivo citaremos como \cite{VIII.D}. \Cf también \cite{VIII.2} para una
exposición sucinta.} que esto significa dos cosas:

\begin{corollaire}[Descenso de homomorfismos de módulos]
\label{VIII.1.2}
\index{descenso de homomorfismos de módulos|hyperpage}%
Sean $g\colon S'\to S$ un morfismo de preesquemas fielmente plano y
cuasicompacto, $F$ y $G$ dos módulos cuasicoherentes sobre~$S$, $F'$ y
$G'$ sus imágenes inversas sobre~$S'$, y, finalmente, $F''$ y $G''$ sus
imágenes inversas sobre~$S''=S'\times_{S}S'$. Consideremos el diagrama de
aplicaciones de conjuntos definido por los funtores de cambio de base
por~$g$, $p_{1}$ y $p_{2}$ (donde $\xymatrix@C=.5cm{p_{1},p_{2}\colon
S'\times_{S}S'\ar@/^2pt/[r]\ar@/_2pt/[r]&S'}$ son las dos proyecciones):
$$
\xymatrix@C=.5cm{\Hom_{S}(F,G)\ar[r] &
\Hom_{S'}(F',G')\ar@/^2pt/[r]\ar@/_2pt/[r] &
\Hom_{S''}(F'',G'').}
$$
Este diagrama es exacto, \ie identifica biyectivamente el primer conjunto
con el igualador de las dos aplicaciones escritas del segundo conjunto
en el tercero.
\end{corollaire}

En otros términos, el funtor de cambio de base
\ifthenelse{\boolean{orig}}
{por~$f$,}
{por~$g$,}
$F\mto F'$, define un funtor \emph{plenamente fiel} de la categoría de
los módulos cuasicoherentes sobre~$S$ en la categoría de los módulos
cuasicoherentes sobre~$S'$ provistos de un dato de descenso relativo
\ifthenelse{\boolean{orig}}
{a~$f$.}
{a~$g$.}
Además:

\begin{corollaire}[Descenso de módulos]
\label{VIII.1.3}
\index{descenso de módulos|hyperpage}%
Para todo módulo cuasicoherente~$F'$ sobre~$S'$, todo dato de descenso
sobre~$F'$ relativo a~$g$ es \emph{efectivo},
\index{dato de descenso efectivo|hyperpage}%
\ie $F'$, con su dato de descenso, es isomorfo a la imagen inversa por~$g$
de un módulo cuasicoherente sobre~$S$ (determinado salvo isomorfismo único
en virtud de~\Ref{VIII.1.2}).
\end{corollaire}

En otros términos, el funtor plenamente fiel anterior es incluso una
\emph{equivalencia}. En la práctica, esto significa que dar un módulo
cuasicoherente sobre~$S$ equivale a dar un módulo cuasicoherente sobre
$S'$ provisto de un dato de descenso relativo a~$g$.

\subsubsection*{Demostración de~\Ref{VIII.1.1}}
Sea
\marginpar{197}
primero $T$ un $S$-preesquema $S$-isomorfo a la suma disjunta de una
familia
\ifthenelse{\boolean{orig}}
{de abiertos}
{de abiertos}
inducidos~$S_{i}$ de~$S$ que recubren~$S$. Es entonces evidente que el
morfismo estructural~$T\to S$ es un morfismo de descenso efectivo
para~$\cal{F}$ (esto significa precisamente que dar un módulo
cuasicoherente~$F$ sobre~$S$ equivale a dar módulos cuasicoherentes~$F_{i}$
sobre los~$S_{i}$ e isomorfismos de pegado~$\varphi_{ji}\colon
F_{i}|S_{i}\cap S_{j}\to F_{j}|S_{i}\cap S_{j}$ que satisfacen la
\ifthenelse{\boolean{orig}}
{conocida condición de cocadenas}
{conocida condición de cociclo}.
En virtud de~VII,~8
se sigue que, para comprobar que $g\colon S'\to S$ es un morfismo de
descenso efectivo para~$\cal{F}$, basta comprobarlo para el morfismo
$g_{T}\colon T'=T\times_{S}S'\to T$ deducido de~$g$ mediante el cambio de
base~$T\to S$. (Obsérvese que la hipótesis sobre~$T\to S$ permanece
estable bajo cualquier cambio de base, de modo que $T\to S$ es, de
hecho, un morfismo \emph{universal} de descenso efectivo para~$\cal{F}$).
Tomando como~$S_{i}$ abiertos afines que recubren~$S$, quedamos reducidos al
caso en que~$S$ es afín.

Entonces $S'$ es reunión finita de abiertos afines y, tomando el
$S$-esquema suma disjunta de estos, se obtiene un $S$-esquema afín~$S_{1}$
y un $S$-morfismo~$S_{1}\to S'$ plano y sobreyectivo. Por tanto, $S_{1}$ es
también fielmente plano sobre~$S$. Así, si se demuestra que un morfismo
afín y fielmente plano es un morfismo de descenso efectivo para~$\cal{F}$
y, por tanto, un morfismo universal de descenso estricto para~$\cal{F}$
(pues la hipótesis es estable bajo cambio de base), se concluye, en
particular, que el morfismo estructural~$S_{1}\to S$ es un morfismo
universal de descenso estricto para~$\cal{F}$; y, como existe un
$S$-morfismo~$S_{1}\to S'$, resultará, por \cite{VIII.D}, que
$g\colon S'\to S$ es, en efecto, un morfismo de descenso estricto
para~$\cal{F}$.

\ifthenelse{\boolean{orig}}{}
{\enlargethispage{.5cm}}%
Quedamos, pues, reducidos al caso en que~$g$ es un morfismo afín y, como
se ha visto, se puede suponer además que~$S$ es afín; por tanto,
\emph{podemos suponer que $S$ y $S'$ son afines}. En este caso,
\Ref{VIII.1.2} equivale al

\begin{lemme}
\label{VIII.1.4}
Sean $A$ un anillo, $A'$ una $A$-álgebra fielmente plana, $M$ y $N$ dos
$A$-módulos, $M'$ y $N'$ los $A'$-módulos obtenidos mediante el cambio
de anillos~$A\to A'$, y $M''$,$N''$ los
$A''=A'\otimes_{A}A'$-módulos obtenidos mediante el cambio de
anillos~$A\to A''$. Entonces la sucesión de aplicaciones de conjuntos
$$
\xymatrix@C=.5cm{\Hom_{A}(M,N)\ar[r] &
\Hom_{A'}(M',N')\ar@/^2pt/[r]\ar@/_2pt/[r] &
\Hom_{A''}(M'',N'')}
$$
es exacta.
\end{lemme}

Como
\marginpar{198}
el homomorfismo~$N\to N'$ es inyectivo (pues $A'$ es fielmente plano
sobre~$A$), se ve que la primera flecha es inyectiva. Queda demostrar que
si un
\ifthenelse{\boolean{orig}}
{$A'$ homomorfismo}
{$A'$-homomorfismo}
$u'\colon M'\to N'$ es compatible con los datos de descenso, entonces
proviene de un $A$-homomorfismo~$u\colon M\to N$. Ahora bien, esto
significa sencillamente que $u'$ aplica el subconjunto~$M$ de~$M'$ en el
subconjunto~$N$ de~$N'$ (la aplicación inducida~$u\colon M\to N$ será
entonces automáticamente $A$-lineal, pues $u'$ es $A'$-lineal, y se ve
igualmente que $u'$ es necesariamente igual a~$u\otimes_{A}A'$). Si
$x\in M$, entonces $u'(x)$ pertenece al igualador del par de aplicaciones
$\xymatrix@C=.5cm{N'\ar@/^2pt/[r]\ar@/_2pt/[r] & N''}$. Quedamos, pues,
reducidos, para demostrar~\Ref{VIII.1.4},
al caso particular siguiente (correspondiente al caso~$M=A$):
